\documentclass[Screen16to9,17pt]{foils}
\usepackage{zencurity-slides}

\selectlanguage{english}

%\externaldocument{unix-audit-security-oevelser}
\externaldocument{\jobname-exercises}

% Start dit home-lab til IT-sikkerhed
% Kun for medlemmer
% Byg dit eget mini-datacenter og lær mens du eksperimenterer.

% Vil du gerne dykke dybere ned i netværk, servere, virtualisering og sikkerhed – og lære det hele i praksis? Så er dette foredrag med efterfølgende øvelsesdag din chance for at komme i gang med dit eget home-lab: et privat testmiljø, hvor du kan eksperimentere med komplekse systemer, uden at noget går galt i produktion. Vi gennemgår, hvordan du designer et fleksibelt og realistisk home-lab – også med et lille budget. Du får konkrete eksempler på opsætninger, værktøjer og projekter, så du hurtigt kan bygge dit eget læringsmiljø derhjemme.

% Vi gennemgår bl.a.:

% • Hvordan du planlægger økonomien: hvad kan købes brugt, og hvad bør du investere i
% • Netværksopsætning og valg af switch/router
% • Servere, mini-pc’er og virtualisering
% • Containere og drift af flere systemer samtidigt
% • Sikkerhed, remote adgang og overvågning
% • Backupstrategi og simple monitoreringsløsninger
% • Idéer til konkrete projekter i dit home-lab
% • Hvornår dit lab bliver så avanceret, at det næsten er et produktionssetup

% Efter foredraget vil du have overblik over, hvordan du samler de første dele til dit eget lab. Og du får inspiration til, hvordan du kan udvide det trin for trin.

% På dagen aftaler vi også en fælles øvelsesdag (online), hvor vi samles og arbejder videre i praksis.

% Nøgleord: TCP/IP, infrastruktur, racks, hardware, servere, virtualisering, containere, sikkerhed

\begin{document}


\mytitlepage{Start Your Home Lab for IT-Security}
{PROSA April 2026}


\hlkprofiluk


\slide{Goals for today}

\hlkimage{4cm}{bornhack-camp-2024.jpg}

\begin{list2}
\item Strategies for Learning starting from Small labs
\item Show various network and server configurations with standard devices
\item What are the costs -- money and time
\item Security for your home lab -- including backups
\item Ideas for home lab projects -- some which I think are important
\end{list2}

Photo is NWWC camp at BornHack 2024, next is BornHack July 15. to 22. 2026 \link{https://bornhack.dk}\\
-- a nice \emph{home lab} with public IP addresses!


\slide{Time schedule}

\hlkimage{4cm}{kame-noanime-small.png}
\begin{list2}
\item 17:00 - 18:15 Introduction and basics
\item 30min Break, talk or spend time with family
\item 18:45 - 19:45 My ideas and strategies
\item 15min small break
\item 20:00 - 21:00 Security, firewalls introduction, how to do larger setups, that mimic production setups
\end{list2}

During the presentation, in breaks etc. please do connect to the network, play with TCP/IP, switches and routers.





\slide{What is a Home Lab}

\begin{quote}
"A homelab is a self-contained IT environment that you design and control. It usually includes basic infrastructure like compute, storage, and networking, along with a hypervisor or virtualization layer. Some setups are modest and run on a single machine, while others resemble mini datacenters with multiple nodes and shared storage such as SAN or NAS."
\end{quote}
Source: StorMagic, "What Is a Homelab? Why IT Pros Are Building Their Own" (stormagic.com, 2025)

\begin{list2}
\item it's your own private sandbox — no stakes
\item no risk to production systems
\item where you can freely experiment with servers, networking, virtualization, security tools,
\item and anything else you're curious about.
\end{list2}




\slide{Prerequisites}

If you are interested in doing experiments -- welcome

If you want to be an expert in IP and network security I recommend doing exercises

\begin{list1}
\item It is recommended to use virtual machines for the exercises
\item Network security and most internet related security work has the following requirements:
\begin{list2}
\item Network experience
\item TCP/IP principles - often in more detail than a common user
\item Programming is an advantage, for automating things -- can be PowerShell, Python, ...
\item Some Linux and Unix knowledge is in my opinion a {\bf necessary skill} for infosec work\\
-- too many new tools to ignore, and lots found at sites like Github and Open Source written for Linux
\end{list2}
\end{list1}




\slide{Course Materials}

This material is in multiple parts:
\begin{list2}
\item Slide show - presentation - this file
\item Additional resources from the internet are linked throughout
\item Wikipedia has a LOT of nice pages about IP protocols, for example:
\end{list2}

\begin{quote}\small
Transport Layer Security (TLS) is a cryptographic protocol designed to provide communications security over a computer network. The protocol is widely used in applications such as email, instant messaging, and voice over IP, but its use in securing HTTPS remains the most publicly visible.
\end{quote}
Source: \url{https://en.wikipedia.org/wiki/Transport_Layer_Security}

\slide{Book: Practical Packet Analysis (PPA)}

We often need a basic understanding of networking!

\hlkimage{4cm}{PracticalPacketAnalysis3E_cover.png}

\emph{Practical Packet Analysis,
Using Wireshark to Solve Real-World Network Problems}
by Chris Sanders, 3rd Edition
April 2017, 368 pp.
ISBN-13:
978-1-59327-802-1
\link{https://nostarch.com/packetanalysis3}

I recommend this book for people new to networking, it has been in HumbleBundle book bundles multiple times


\slide{Teaching and Learning}

Strategies for Learning
\begin{list1}
\item Typically a home lab, as the name says, is for learning at home
\item You want to learn something!
\item Learning is different for different people -- I learn great through books and experiments
\vskip 2cm
\item Learning something has different parts
\begin{list2}
\item Should be \emph{safe} -- experiment freely without risking production systems
\item Hands-on practice builds real-world confidence and competence
\item Master the basics first -- avoid unnecessary frustration by building on solid foundations
\end{list2}
\end{list1}


\slide{Hiim and Hippe's Didactical relationship Model}
\hlkimage{8cm}{images/hiim-hippe-relationsmodel.png}
Source: \emph{Læring gennem oplevelse, forståelse og handling} 2007 Hiim, Hilde \& Hippe, Else

\begin{list2}
\item Learning Conditions (Prerequisites), Settings (Framework Conditions), Learning Goals and Objectives
\item Content (Curriculum and Subject Matter), Learning Process, Assessment (Evaluation of Learning)
\end{list2}

\slide{Bloom's Taxonomy}

\hlkimage{10cm}{images/bloom.png}
Source: Anderson, L., \& Krathwohl, D. R. (2001). \emph{A taxonomy for learning, teaching and assessing: A revision of Bloom's Taxonomy of Educational Objectives.} New York: Longman. -- revised 2001 version

The levels build upon each other -- a learner must master lower levels  before progressing to higher-order thinking.

%Bloom’s is used in three ways: to set learning outcomes, to structure classroom activities and to assess students’ progress.
%Source: Didactic Plan for IT Security PBA KEA.dk -- now EK.dk



\slide{Setting Goals }

\begin{quote}
  “A goal without a plan is just a wish.”\\
  ― Antoine de Saint-Exupéry
\end{quote}

\hlkimage{8cm}{pawel-janiak-dxFi8Ea670E-unsplash.jpg}

\begin{list1}
\item Point you towards resources, so you can get started
\item List a few core concepts I think you should know and learn
\hskip 2cm {\footnotesize Photo by Pawel Janiak on Unsplash}
\end{list1}


\slide{Core Concepts}

Information Security is a huge domain:

\begin{quote}
The (ISC)² CBK is a collection of topics relevant to cybersecurity professionals around the world. It establishes a common framework of information security terms and principles which enables cybersecurity and IT/ICT professionals worldwide to discuss, debate and resolve matters pertaining to the profession with a common understanding, taxonomy and lexicon.
\end{quote}
Source: \link{https://www.isc2.org/Certifications/CBK}

List of 8 domains in CISSP CBK: Security and Risk Management, Asset Security,
Security Architecture and Engineering, Communications and Network Security, Identity and Access Management, Security Assessment and Testing, Security Operations, Software Development Security

- then add all the news about new tools, exploits, and networking


\slide{Learning Nmap at different levels}

\hlkimage{10cm}{trinity-nmapscreen-hd-cropscale-418x250.jpg}

To illustrate this, I will use the example of:\\
Nmap - a very famous port scanner.

Unfortunately there are about 100 options, and the man page is some 3100 lines ...

\slide{Prerequisite knowledge}

Plan: You want to learn Nmap!

{\large Often we combine our knowledge with skills into competence, \\
which enable us to perform some job, task or function.}

\begin{list2}
\item Knowledge level: What is a port scanner\\
Need to know TCP/IP, IP address, ports and services -- example HTTP 80/tcp, TCP session setup
\end{list2}

So get this sorted out first, otherwise you cannot understand what Nmap does, and output returned

\slide{Skills are needed}

\hlkimage{9cm}{nmap-zenmap.png}


\begin{list2}
\item Skills level: Running a port scanner\\
Need to have operating system -- luckily Nmap supports Mac, Windows, Linux, ...
\item My recommendation: create a virtual machine with Kali Linux
\end{list2}

\slide{Combined it becomes a Competence}

\begin{alltt}\footnotesize
full-tcp-scan: nmap -p 1-65535 -A -oA full-tcp-scan -iL targets
dns-scan: nmap -sU -p 53 --script=dns-recursion -iL targets -oA dns-recursive
bgpscan: nmap -A -p 179 -oA bgpscan -iL \$LINKNET
dns-recursive: nmap -sU -p 53 --script=dns-recursion -iL targets -oA dns-recursive
php-scan: nmap -sV --script=http-php-version -p80,443 -oA php-scan -iL targets
scan-vtep-tcp: nmap -A -p 1-65535 -oA scan-vtep-tcp 192.0.2.77 192.0.2.78
snmp-10.x.y.0.gnmap: nmap -sV -A -p 161 -sU --script=snmp-info -oA snmp-10xy 10.x.y.0/19
snmpscan: nmap -sU -p 161 -oA snmpscan --script=snmp-interfaces -iL targets
sshscan: nmap -A -p 22 -oA sshscan -iL targets
vncscan: nmap -A -p 5900-5905 -oA vncscan -iL targets
\end{alltt}

\begin{list2}
\item Competence level: Running a quality port scan of an enterprise\\
Need to have plan for scanning, know which scan functions to use
\end{list2}

My recommendation: work through a 4 hour course with Nmap as the subject


\slide{OSI Model and Internet Protocols}

\hlkimage{10cm,angle=90}{images/compare-osi-ip.pdf}

\slide{Recommended technologies to learn}


So to accomplish the goal of using Nmap efficiently you need some basics

Networking: Basic Protocols from the Internet Protocols suite IP/TCP, or TCP/IP
\begin{list2}
\item Network Layer: Ethernet, Address Resolution Protocol (ARP), IPv4 and ICMP\\
Later add Wi-Fi and IPv6
\item Transport Layer: Transmission Control Protocol (TCP) and User Datagram Protocol (UDP)
\item Common upper layer: Dynamic Host Configuration Protocol (DHCP), Domain Name System (DNS),
Hypertext Transfer Protocol (HTTP)\\
Later add the encrypted/secure versions like Hypertext Transfer Protocol Secure (HTTPS) which uses Transport Layer Security (TLS)
\end{list2}

Pro tip: always say Ethernet frames and IP packets. No one uses datagram anymore.

Pro tip: If you \emph{really know DNS} you can make a huge impact in the malware area!

\slide{Books and courses}

%\hlkimage{}{}

How:

I like to learn new concepts from books
\begin{list2}
\item Have a clear structure, less confusion
\item They go from a basic level towards a complete goal
\item Often have exercises available with nice progression
\item Often you can get Humble bundles with many books for \$25
\item Some books are "free" if you give your email address, example\\
\emph{Web Application Security}, Andrew Hoffman, 2020, ISBN: 9781492053118 - was available for download for free through Nginx at some point
\end{list2}

Pro Tip: all my courses and exercise booklets are available on Codeberg!

\slide{ANSM Cycle}

\hlkimage{7cm}{NSM_Phases-300x238.png}

\begin{list1}
\item ANSM chapter 1: The Practice of Applied
Network Security Monitoring
\begin{list2}
\item Vulnerability-Centric vs. Threat-Centric Defense
\item The NSM cycle: collection, detection, and analysis
\item Full Content Data, Session Data, Statistical Data, Packet String Data, and Alert Data
\item Security Onion is nice, but a bit over the top - quickly gets overloaded
\end{list2}
\end{list1}

Sooo much in the first chapter!

\slide{Baseline Skills}

This book is great, since it lays ouy the different subjects nicely:
\begin{list2}\small
\item Threat-Centric Security, NSM, and the NSM Cycle
\item TCP/IP Protocols
\item Common Application Layer Protocols
\item Packet Analysis
\item Windows Architecture
\item Linux Architecture
\item Basic Data Parsing (BASH, Grep, SED, AWK, etc)
\item IDS Usage (Snort, Suricata, etc.)
\item Indicators of Compromise and IDS Signature Tuning
\item Open Source Intelligence Gathering
\item Basic Analytic Diagnostic Methods
\item Basic Malware Analysis
\end{list2}

Source: \emph{Applied Network Security Monitoring Collection, Detection, and Analysis}, Chris Sanders \& Jason Smith, 2014


\slide{Course Network / Small Networks}
.
\hlkrightpic{65mm}{-1cm}{sample-network.png}

\begin{list1}
\item I often have a course network with me which has the following information:
\begin{list2}
\item Routers -- computer with OpenBSD and portable routers
\item Switches Juniper and TP-Link
\item Wireless access-points Unifi AP
% \item IPv4 addresses: 151.216.32.0/21 total 2048
% \item IPv6 addresses route6: 2001:678:9ec::/48
% \item Autonomous system number: AS208647 BornHack\\
\url{https://en.wikipedia.org/wiki/Autonomous_system_(Internet)}
\end{list2}
\end{list1}

\vskip 2cm
{\bf I will highly recommend that you have a production network and a testing network!}\\
If you get toooo frustrated with the home lab -- go relax, listen to music, stream a movie etc. on the production network \smiley




\slide{Testing network the legal issues}

{\bfseries Straffelovens paragraf 263 Stk. 2. Med bøde eller fængsel indtil 1 år og 6 måneder straffes den, der uberettiget skaffer sig adgang til en andens oplysninger eller programmer, der er bestemt til at bruges i et informationssystem. }

\begin{list2}
\item Danish law about hacking
\item Please check with your legal department, or be careful
\item We {\bf always} contact network between us and the network to be tested
\item Be good netizens
\end{list2}




\slide{Hackerlab setup -- easy}

\hlkimage{8cm}{hacklab-1.png}

\begin{list2}
\item I recommend getting a hackerlab running on your laptop
\item Hardware: modern laptop which has CPU virtualization\\
Dont forget to check BIOS settings for virtualization
\item Software: your favorite OS: Windows, Mac, Linux
\item Virtualization software: VMware, Virtual box, HyperV
\item Hacker software: Kali as a Virtual Machine \link{https://www.kali.org/}
%\item Soft targets: Metasploitable, Windows 2000, Windows XP, ...
\end{list2}


\slide{Internet in a Box}

\hlkrightimage{10cm}{internet-in-a-box.jpeg}

The main purpose of showing this box, are:
\begin{list2}
\item These are standard devices, Juniper EX3300 cheap oldish, works great
\item Managed switches are a must! You can learn by buying cheap ones,\\
like the TP-Link T1500G-10PS  shown, VLAN, SNMP, Syslog ...
\item Multiple systems created using PC Engines APU2C4 (really D4)\\
running OpenBSD, Unbound, Suricata, Zeek, DHCP, \\
router advertisement, PF firewall - explicit and nice ICMPv6 filtering ...
\item Attack systems compact PCs or laptop
\item Creating a home lab is not expensive, \\
bought the Arista 7150 24-port 10G used on ebay
\end{list2}

You should have similar (or better) devices in your production network, and they can be
configured to do a LOT more than you use them for right now



\slide{Operating Systems}


\slide{Kali Linux the new backtrack}

\hlkimage{16cm}{kali-linux.png}

\begin{list1}
\item Kali \link{http://www.kali.org/}
\item Use the documentation on Kali website for installation
\item Use tools web sites for documentation about tools, like \link{https://nmap.org/} and \link{http://www.hping.org/}
\end{list1}


\slide{Debian with XFCE}


\slide{Virtualisation}

Use it,

create many virtual machines:


\begin{list2}
\item Switch context easily, one project - one (or more) VMs
\item Create a new VM -- keep one updated that you can clone
\item Install crap -- delete a VM, retry
\end{list2}


\slide{Demo clone a VM }


Show it on Qubes OS

Show the KVM virt-manager shortly on Bob


\slide{various network and server configurations}

\slide{Minimum Lab -- recommendation}

Buy a router -- can be a physical router Opal

\slide{What are the costs -- money and time}

\slide{Buying stuff}

Try to buy stuff that will last, can be used for more things

- you probably want to expand the lab, do more complex stuff later


\slide{Opal router }

slide from basic tcp and security


\slide{A small network home lab}

A switch, can be a single uplink port on the production network

A router physical or virtualized

A small \emph{server} -- can be a mini PC

\slide{Demo }

My regular course network

\begin{list2}
\item Switch -- managed
\item Physical router or VM with software
\item Server -- mini-PC for virtualized servers
\item Cables and a power strip
\end{list2}

Budget < 5.000DKK

\vskip 2cm
AP and wireless left out, but a physical router might have it anyway

\begin{list2}
\item
\item
\item
\item
\end{list2}

\slide{Buying Used and Refurbished}

ebay

refurb devices -- my PCs

the old pile of stuff in the office!

\slide{My ideas and strategies}



\slide{Security}

Security, firewalls introduction, how to do larger setups, that mimic production setups

\slide{Security for your home lab -- including backups}


\slide{Ideas for home lab projects -- some which I think are important}


\slide{Conclusion}

\begin{list1}
\item You really should try doing experiments
\item Investigate your existing devices -- you may be able to re-use devices
\item Choose which devices does which part -- production or testing
\item Plan ahead -- both which things you want to test, but how to make it efficient
\end{list1}

\myquestionspage


\end{document}
